\sgaStableTarget{sga6.e05.eq.3-5-1}@@M0@@

If @@M1@@ is a @@M2@@-ring, then for all @@M3@@ we have:
\sgaStableTarget{sga6.e05.eq.3-5-2}@@M4@@

We denote by @@M7@@ the @@M5@@-th universal polynomial of the law @@M6@@; it is a polynomial in @@M8@@ whose isobaric components have weights ranging from @@M9@@ to @@M10@@ with respect to all the variables @@M11@@, (1.15) and which is symmetric with respect to these two groups of variables. For @@M12@@, we therefore have:
\sgaStableTarget{sga6.e05.eq.3-5-3}@@M13@@

Finally, note that the law @@M14@@ has a unit element, namely @@M15@@, that is, @@M16@@; it will sometimes be denoted by @@M17@@.

\sgaStableTarget{sga6.e05.para.3-6}3.6. The ``operations @@M18@@'' are defined by the commutative diagram:
\sgaStableTarget{sga6.e05.eq.3-6-1}@@M19@@
i.e. by the composite morphism of functors:
@@M20@@
and will be denoted by @@M21@@. They are defined by their value on @@M22@@; we have:
